\chapter{Osservabilit\`a}
\label{cap:21-osservabilita}

In una piattaforma pensata come strumento di misura, l'osservabilit\`a non \`e un
corredo operativo: \`e l'apparato sperimentale. Ogni risultato riportato nella
Parte V \`e stato ottenuto leggendo le metriche descritte in questo capitolo, e
diverse decisioni di progetto della Parte III sono state prese sulla base di
grandezze che qui sono definite.

\section{Il modello}

Le metriche sono Micrometer esportate in formato Prometheus~\cite{prometheus} su
\code{/actuator/prometheus}, porta \code{8081}. La separazione dalla porta
dell'API (\cref{cap:04-architettura-control-plane}) fa s\`i che lo scraping non
condivida code n\'e limiti di ammissione con il traffico misurato.

La quasi totalit\`a delle metriche \`e etichettata per funzione. Il ciclo di vita
dei meter segue quello della funzione: creati alla registrazione, rimossi alla
cancellazione (\cref{cap:06-nucleo}).

\section{Il catalogo}

\subsection{Nucleo}

\begin{center}
\small
\begin{adjustbox}{max width=\textwidth}
\begin{tabular}{@{}llp{6.2cm}@{}}
\toprule
\textbf{Metrica} & \textbf{Tipo} & \textbf{Significato} \\
\midrule
\code{function\_enqueue\_total} & counter & Accodamenti. \\
\code{function\_dispatch\_total} & counter & Dispatch verso il runtime.
Sorgente della metrica \code{rps} dell'autoscaler. \\
\code{function\_success\_total} & counter & Esiti positivi. \\
\code{function\_error\_total} & counter & Esiti negativi. \\
\code{function\_retry\_total} & counter & Riprove. \\
\code{function\_latency\_ms} & timer & \textbf{Tempo di servizio}: da dispatch a
completamento. \\
\code{function\_e2e\_latency\_ms} & timer & Tempo di sojourn: da ammissione a
completamento. \\
\code{function\_queue\_wait\_ms} & timer & Attesa in coda. \\
\code{function\_cold\_start\_ms} & timer & Durata dell'inizializzazione. \\
\code{function\_init\_duration\_ms} & timer & Come sopra, dal runtime. \\
\bottomrule
\end{tabular}
\end{adjustbox}
\end{center}

La distinzione fra le prime tre metriche temporali \`e la pi\`u importante
dell'intero capitolo, ed \`e il fondamento del \cref{cap:12-concurrency-control}:

\[
  \underbrace{\code{function\_e2e\_latency\_ms}}_{\text{sojourn}}
  \;=\;
  \underbrace{\code{function\_queue\_wait\_ms}}_{\text{attesa}}
  \;+\;
  \underbrace{\code{function\_latency\_ms}}_{\text{servizio}} .
\]

Che le tre esistano separatamente \`e ci\`o che ha reso possibile stabilire che
la seconda \`e sei volte la terza sotto carico, e quindi che un controllore che
decide dalla terza sta ottimizzando la grandezza sbagliata.

\subsection{Code}

\begin{center}
\small
\begin{adjustbox}{max width=\textwidth}
\begin{tabular}{@{}llp{6.2cm}@{}}
\toprule
\textbf{Metrica} & \textbf{Tipo} & \textbf{Fornita da} \\
\midrule
\code{function\_queue\_depth} & gauge & \code{async-queue} \\
\code{function\_inFlight} & gauge & \code{async-queue} \\
\code{function\_effective\_concurrency} & gauge & \code{async-queue} \\
\code{sync\_queue\_depth} & gauge & \code{sync-queue}, globale e per funzione \\
\code{sync\_queue\_wait\_seconds} & gauge & \code{sync-queue}, attesa stimata \\
\code{sync\_queue\_admitted\_total} & counter & \code{sync-queue} \\
\code{sync\_queue\_rejected\_total} & counter & \code{sync-queue} \\
\code{sync\_queue\_timedout\_total} & counter & \code{sync-queue} \\
\bottomrule
\end{tabular}
\end{adjustbox}
\end{center}

\subsection{Politiche}

\begin{center}
\small
\begin{adjustbox}{max width=\textwidth}
\begin{tabular}{@{}llp{6cm}@{}}
\toprule
\textbf{Metrica} & \textbf{Tipo} & \textbf{Fornita da} \\
\midrule
\code{function\_target\_inflight\_per\_pod} & gauge & \code{concurrency-control} \\
\code{function\_concurrency\_controller\_mode} & gauge &
\code{concurrency-control} \\
\code{nanofaas\_offload\_total} & counter & \code{offload}, con etichetta
\code{trigger} \\
\code{nanofaas\_offload\_failure\_total} & counter & \code{offload} \\
\code{controlplane\_runtime\_config\_revision} & gauge & \code{runtime-config} \\
\bottomrule
\end{tabular}
\end{adjustbox}
\end{center}

Nessuna di queste alimenta una decisione: sono tutte \emph{descrittive della
decisione}. \`E una categoria che vale la pena isolare concettualmente, perch\'e
risponde a una domanda che le metriche di stato non possono risolvere: dato un
comportamento osservato, quale politica era attiva e che cosa aveva deciso.

\section{I due profili}

\begin{lstlisting}[language=Java]
private static final Set<String> ADVANCED_METRICS = Set.of(
        "function_cold_start_total", "function_warm_start_total",
        "function_effective_concurrency", "function_target_inflight_per_pod",
        "function_concurrency_controller_mode", "gateway_service_target_load",
        "controlplane_runtime_config_revision", ...);
\end{lstlisting}

Il profilo \code{basic} (default) esporta il nucleo. Il profilo
\code{advanced} aggiunge le metriche descrittive delle politiche e --- pi\`u
significativamente --- attiva gli istogrammi per percentili sui quattro timer
per funzione.

Il costo che i due profili arbitrano \`e quello degli istogrammi. Un timer
Micrometer senza istogramma mantiene conteggio, somma e massimo: tre valori. Con
gli istogrammi mantiene una serie di secchi, e ogni secchio \`e una serie
temporale in Prometheus. Per una piattaforma con molte funzioni la differenza
nella cardinalit\`a esportata \`e sostanziale.

La conseguenza pratica \`e che i percentili --- p95, p99 --- sono disponibili
soltanto sotto \code{advanced}. Poich\'e i risultati della Parte V sono espressi
prevalentemente in percentili, tutte le campagne sperimentali girano su quel
profilo.

Il profilo \`e validato all'avvio: un valore diverso da \code{basic} o
\code{advanced} fa fallire il contesto con un messaggio esplicito, anzich\'e
ripiegare in silenzio su un default.

\section{Le letture canoniche}

La documentazione operativa del progetto raccoglie le interrogazioni PromQL di
riferimento. Le pi\`u istruttive sono quelle che \emph{combinano} due metriche,
perch\'e ciascuna combinazione codifica una diagnosi.

\begin{lstlisting}[language=bash, caption={Tasso di errore per funzione.}]
rate(function_error_total{function="word-stats"}[5m])
  / (rate(function_success_total{function="word-stats"}[5m])
     + rate(function_error_total{function="word-stats"}[5m])) * 100
\end{lstlisting}

\begin{lstlisting}[language=bash, caption={Pressione di ammissione sincrona.}]
rate(sync_queue_rejected_total{function="word-stats"}[5m])
  / rate(sync_queue_admitted_total{function="word-stats"}[5m])
\end{lstlisting}

\begin{lstlisting}[language=bash, caption={Carico effettivo per replica: il
segnale su cui l'autoscaler decide.}]
rate(function_dispatch_total{function="word-stats"}[1m])
  / max by (function) (kube_deployment_status_replicas{deployment="fn-word-stats"})
\end{lstlisting}

\subsection{La guida alla lettura della contesa}

Il progetto documenta quattro regole di interpretazione che meritano di essere
riportate, perch\'e sono il tipo di conoscenza che di norma resta implicita:

\begin{itemize}[leftmargin=*]
  \item \code{sync\_queue\_depth} crescente con \code{function\_dispatch\_total}
  piatto significa che l'ammissione riesce pi\`u in fretta di quanto gli slot si
  riaprano.
  \item \code{sync\_queue\_rejected\_total} alto con profondit\`a bassa indica un
  rifiuto per attesa stimata, non un esaurimento di capacit\`a.
  \item \code{function\_dispatch\_total} che cresce mentre successi ed errori
  restano indietro indica latenza di completamento, non iniquit\`a dello
  scheduler.
  \item Profondit\`a persistente con poche richieste in volo indica una funzione
  limitata dagli slot; profondit\`a persistente con molte richieste in volo
  indica un runtime lento.
\end{itemize}

L'ultima coppia \`e la pi\`u utile e non \`e ovvia: la stessa profondit\`a di coda
ha due cause opposte, distinguibili solo confrontandola con il contatore delle
richieste in volo.

\section{Il tasso di errore \`e ambiguo, e il progetto lo disambigua}

Un problema ricorrente nel misurare piattaforme FaaS: il tasso di errore
mescola i fallimenti della piattaforma con le risposte non-2xx deliberate delle
funzioni. Una funzione che restituisce correttamente \code{404} per una risorsa
inesistente non \`e un errore della piattaforma, ma da fuori \`e indistinguibile.

Il meccanismo della busta di risposta (\cref{cap:17-runtime-funzione}) risolve
l'ambiguit\`a alla radice: una risposta decisa dall'handler porta
\code{X-NanoFaaS-Function-Status: true} ed \`e contabilizzata come successo dalla
piattaforma, qualunque sia il suo codice di stato. \code{function\_error\_total}
conta quindi soltanto i fallimenti che la piattaforma si attribuisce.

\`E una precondizione per il gate sperimentale del
\cref{cap:23-prestazioni-base}, che richiede tasso di errore esattamente zero:
un tale gate sarebbe inapplicabile se il contatore includesse le risposte
applicative negative.

\section{Log e tracing}

I log sono strutturati con cinque campi: \code{executionId},
\code{functionName}, \code{traceId}, \code{attempt}, \code{status}. Il tracing
\`e la propagazione di \code{X-Trace-Id} lungo l'intera catena, con supporto per
i contesti W3C (\code{traceparent}, \code{tracestate}) sul percorso di offload.

Una regola \`e enunciata esplicitamente e vale la pena riportarla:

\begin{quote}\itshape
Il runtime non sintetizza identificatori di esecuzione segnaposto per il
logging. Se n\'e l'header della richiesta n\'e \code{EXECUTION\_ID} sono
disponibili, il contesto diagnostico resta vuoto per \code{executionId}.
\end{quote}

Un identificatore inventato riempirebbe il campo e renderebbe la correlazione
\emph{apparentemente} possibile ma sbagliata. Un campo vuoto \`e onesto.

\subsection{La consegna delle callback \`e limitata e pu\`o essere scartata}

Un dettaglio con conseguenze sperimentali. La consegna della callback di
completamento dal runtime Java \`e asincrona e limitata: \code{/invoke}
restituisce senza attenderla, e quando il dispatcher \`e saturo la callback viene
\emph{scartata} con un avviso. Il runtime Go si comporta allo stesso modo e
incrementa una metrica di callback scartate.

Una callback scartata produce un'esecuzione che non raggiunge mai uno stato
terminale --- esattamente il caso che il terzo orizzonte di eviction dello
\code{ExecutionStore} deve ripulire (\cref{cap:06-nucleo}). Che il runtime Go
esponga un contatore dedicato e quello Java soltanto un avviso nel log \`e
un'asimmetria che rende il fenomeno misurabile su un SDK e non sull'altro.

Le riprove della callback sono limitate ai fallimenti ritentabili: errori di
trasporto, \code{408}, \code{429} e \code{5xx}. Gli altri \code{4xx} sono
considerati permanenti. \`E la distinzione corretta --- un \code{400} sulla
callback significa che il control plane ha rifiutato il corpo, e ripeterlo non
lo far\`a accettare.

\section{Salute}

Due sonde: \code{/actuator/health/liveness} e
\code{/actuator/health/readiness}, entrambe sulla porta di management. Sono la
base delle \code{livenessProbe} e \code{readinessProbe} del Deployment del
control plane.

\section{Copertura contro le regressioni}

Il progetto adotta una posizione metodologica che vale la pena riportare per
esteso:

\begin{quote}\itshape
Il repository include test strutturali di regressione sul percorso caldo anzich\'e
microbenchmark assoluti. Asseriscono progresso e comportamenti sensibili
all'allocazione, come il riuso del replay, il forward progress della coda
sincrona dietro una testa bloccata, e l'equit\`a asincrona fra funzioni calde e
fredde. [\ldots] Nel tarare il comportamento delle code, si preferisca
preservare quelle garanzie strutturali anzich\'e inseguire un numero fisso di
temporizzazione locale. Le temporizzazioni assolute sono sensibili all'ambiente,
le garanzie di equit\`a e riuso non lo sono.
\end{quote}

\`E la risposta corretta a un problema noto dei test di prestazione: un test che
asserisce ``questa operazione impiega meno di $X$ millisecondi'' fallisce sulla
macchina di un collega, e viene disattivato. Un test che asserisce ``dietro una
testa bloccata, un'altra funzione riesce comunque a progredire'' cattura la
propriet\`a che interessava, ed \`e vero su qualunque macchina.
